\uuid{RLPV}
\titre{Deux dés à trois faces}
\auteur{worms}
\organisation{exo7}
\licence{CC BY-NC-SA 4.0}
\datecreate{2017-08-28}
\niveau{L2}
\module{Probabilité et statistique}
\chapitre{Probabilité discrète}
\sousChapitre{Probabilité et dénombrement}
\difficulte{}
\theme{Probabilités, événements}
\exo7qcmid{worms.worms.014}
\serie{ordonnee}

\contenu{
\texte{
Une expérience consiste à lancer deux dés à 3 "faces" (si, si, ça existe ! \'Equiprobables bien entendu.). On note $A_i=$"le premier dé vaut $i$" et $B_i=$"le second dé vaut $i$" pour chaque $i\in\{1,2,3\}$, ainsi que $S_k=$"la somme des deux dés vaut au plus $k$" pour $k\in \{2,3,4,5,6\}$. On note $\Omega$ l'univers associé à cette expérience.
}

\begin{question}
\qclassification{260.01}
Parmi les descriptions ci-dessous, laquelle/lesquelles désigne(nt) une partition de $\Omega$ ?
\begin{answers}
\good{$\{B_1,B_2,B_3\}$}
\good{$\{A_1\cup A_2,A_3\}$}
\bad{$\{S_2,S_3,S_4,S_5,S_6\}$}
\end{answers}
\begin{explanations}
$\{S_2,S_3,S_4,S_5,S_6\}$ n'est pas du tout une partition car les événements $S_k$ ne sont pas incompatibles deux à deux.
\end{explanations}
\end{question}

\begin{question}
\qclassification{260.01}
Parmi les affirmations suivantes, laquelle/lesquelles est/sont erronée(s) ou n'a/n'ont aucun sens ?
\begin{answers}
\good{$(A_1\cup A_2\cup A_3)^c=B_1\cap B_2$}
\bad{$\mathbb{P}(\Omega)=\cup_{i=1}^3\mathbb{P}(A_i)$}
\good{$(A_1\cup A_2)\cap(B_1\cup A_3)=(A_1\cap B_1)\cup (A_2\cap B_1)$}
\bad{$S_3=(A_1\cap B_1)\cup(A_1\cap B_2)$}
\bad{$S_5^c=A_3\cup B_3$}
\good{$\mathbb{P}(B_j|A_i)=\mathbb{P}(B_j)$ \ $(\forall (i,j)\in \{1,2,3\}^2)$}
\end{answers}
\begin{explanations}
Pour le $(A_1\cup A_2\cup A_3)^c$,  on vérifie que les deux membres sont vides. Pour $\mathbb{P}(\Omega)=\ldots$, on ne réunit pas des probabilités ! Pour $(A_1\cup A_2)\cap\ldots$ on le vérifie en développant à gauche et en rencontrant certaines intersections vides.   Pour $S_3=\ldots$ il manque $A_2\cap B_1$. Pour $S^c_5=\ldots$ c'est une intersection plutôt. Pour $\mathbb{P}(B_j|A_i)=\ldots$ chaque $B_j$ est indépendant de chaque $A_i$, oui.
\end{explanations}
\end{question}

\begin{question}
\qclassification{260.01}
Que vaut la probabilité de $S_5$ ?
\begin{answers}
\bad{$5/6$}
\bad{$7/9$}
\good{$8/9$}
\end{answers}
\begin{explanations}
Comme $S_5^c = A_3 \cap B_3$ et que $A_3$ et $B_3$ sont indépendants de probabilités $1/3$ alors
on a $\mathbb{P}(S_5) = 1-1/9 = 8/9$.
\end{explanations}
\end{question}

\begin{question}
\qclassification{260.01}
Que vaut la probabilité de $S_5\backslash S_2$ ?
\begin{answers}
\bad{$2/3$}
\good{$7/9$}
\bad{$6/9$}
\bad{L'énoncé ne veut rien dire}
\end{answers}
\begin{explanations}
$S_5\backslash S_2$ signifiant $S_5$ privé de $S_2$, et comme $S_2$ est inclus dans $S_5$, alors
$\mathbb{P}(S_5\backslash S_2) = \mathbb{P}(S_5) - \mathbb{P}(S_2)
= 8/9-1/9 = 7/9$.
\end{explanations}
\end{question}

\begin{question}
\qclassification{260.01}
Que vaut la probabilité conditionnelle de $S_5$ sachant $S_2$ ?
\begin{answers}
\bad{$1/8$}
\good{$1$}
\bad{$8/9$}
\bad{L'énoncé ne veut rien dire}
\end{answers}
\begin{explanations}
C'est $1$ car $S_2$ est inclus dans $S_5$.
\end{explanations}
\end{question}
}
